\section{Open questions}

Five concrete follow-up questions from a genuine re-read of the paper against
this replication. Each item lists (i) the question, (ii) the basis (what in
the paper or replication motivates it), and (iii) concrete next-step probes
that a follow-on replicator or the original authors could execute.

\begin{enumerate}

\item \textbf{Gamma vs alternative distributions --- is the R$^2 > 0.999$ claim
robust to family selection?} \\
\emph{Basis:} The paper reports a Gamma fit with $R^2 > 0.999$ across all 16
monoenergetic beams but does not compare against log-normal, negative-binomial,
truncated-exponential, or Gamma-mixture alternatives. Physical intuition
suggests high-LET tracks (dense track-core + sparse penumbra) could produce a
bimodal complexity distribution, for which a Gamma-mixture would be more
natural. Raw SDD histograms would answer this at zero extra cost, but they
are not deposited. \\
\emph{Next steps:} (a) request raw SDD per-event damage tables from
Bertolet/Schuemann or barring that, the per-beam damage-count histograms;
(b) fit each of \{Gamma, log-normal, negative-binomial, Gamma-mixture-of-2,
truncated-exponential\} by MLE and tabulate $\Delta$AIC/$\Delta$BIC per beam;
(c) run Hartigan's dip test for bimodality on the high-LET beams
($y_F > 100$ keV/$\mu$m).

\item \textbf{Is quadratic-in-$y_F$ the right functional form for
$\alpha(y_F), \beta(y_F)$, or is the paper using a convenience polynomial
where a saturating physical form belongs?} \\
\emph{Basis:} The paper fits $\alpha, \beta$ as quadratics in $y_F$ with no
mechanistic justification, and $\beta(y_F)$ crosses zero at $y_F \approx 175$
keV/$\mu$m --- inside the paper's own stated validity window. The paper uses
linear-plus-saturating-exp forms for other damage-count observables (SBI, BDI),
so it is asymmetric to switch to unrestricted polynomials only for the Gamma
parameters. \\
\emph{Next steps:} (a) once raw histograms are obtained, re-fit
$\alpha(y_F), \beta(y_F)$ with \{quadratic, linear+saturating-exp,
Michaelis-Menten, power-law\} and compare AIC; (b) enforce $\beta > 0$
(e.g., fit $\log \beta$ or a saturating form) and check extrapolation to
$y_F \sim 300$ keV/$\mu$m (relevant for BNCT alphas); (c) publish a code-note
documenting the current $y_F \approx 175$ keV/$\mu$m ceiling.

\item \textbf{Does the Gamma-complexity distribution predict cell survival
better under a mechanistic repair model than under the paper's sigmoid?} \\
\emph{Basis:} The paper explicitly describes its sigmoid repair model as
``qualitative'', which weakens the downstream RBE-vs-LET curves (Figure 5)
as a test of the Gamma-complexity hypothesis. A mechanistic model (kinetic
NHEJ/HR, LEM-IV, or Medras) would be a stronger downstream test. \\
\emph{Next steps:} (a) take the reproduced Gamma complexity distributions and
feed them into McMahon et al.'s Medras kinetic repair model (open source,
\texttt{github.com/sjmcmahon/Medras-MC}); (b) compare Medras+Gamma-complexity
vs sigmoid+Gamma-complexity predictions against the same V79-4 and irs1
experimental survival data (Hill 2004, Jones 1987); (c) if Medras+Gamma-complexity
outperforms the sigmoid variant, that strengthens the case that the
Gamma-complexity distribution is capturing physically-relevant damage
information rather than a coincidence of the fit.

\item \textbf{How sensitive are $(\alpha, \beta)$ to scoring-volume geometry?} \\
\emph{Basis:} The paper fixes the scoring volume as a 1\,$\mu$m-diameter
sphere. But complexity is inherently scale-dependent: nucleosome-scale
(100 nm) volumes pick up track-core physics, whereas nucleus-scale (5+ $\mu$m)
volumes smear over independent tracks. No sensitivity study across scoring
volumes is reported. \\
\emph{Next steps:} (a) re-run TOPAS-nBio for one reference beam
(4-MeV alpha, $y_F \sim 95$ keV/$\mu$m) with three scoring-volume diameters:
100 nm, 1\,$\mu$m (baseline), 5\,$\mu$m; (b) fit the resulting complexity
histograms with a Gamma and record $R^2$ and $(\alpha, \beta)$ per scale;
(c) if the Gamma family remains $R^2 > 0.99$ across all three scales but
$(\alpha, \beta)$ shift systematically, publish a scale-transformation law;
if the Gamma family breaks down at one extreme, that bounds the physical
regime of the Gamma model.

\item \textbf{Can the Gamma-complexity model be validated directly against
plasmid DSB-yield experiments (repair-free)?} \\
\emph{Basis:} Plasmid irradiation experiments produce clean SSB/DSB yields
per Gy per Da of DNA, without any repair-model confound. Milligan, Ward,
Belli, Fulford and others published such data at multiple LETs. The paper
tests only against cell-survival data (Figure 5), which convolves the
complexity distribution with an unknown repair sigmoid. Plasmid data would
be a much cleaner test of the physics layer alone. \\
\emph{Next steps:} (a) curate published plasmid DSB-yield vs LET
measurements (Milligan JR et al., Radiat Res 1993; Fulford et al.,
Int J Radiat Biol 2001; Belli et al. plasmid series); (b) for each
plasmid experiment's beam parameters, compute $y_F$ via standard
microdosimetry and use MGM to predict $N_{\text{sites,DSB}}$ per track and
the complex-to-simple DSB ratio; (c) compare predicted vs measured yields ---
if MGM reproduces the plasmid data across LET, that is a strong, repair-free
validation of the physics layer; if not, the failure is isolated to the
biological repair layer.

\end{enumerate}
